\section{Related Work}
\label{sec:related}

\paragraph{SPD geometry and matrix completion.}
The affine-invariant geometry of positive matrices is classical
\citep{bhatia2007positive,pennec2006riemannian}.  Positive definite and
covariance completion usually infer missing matrix entries from sparsity or
graph constraints \citep{dempster1972covariance,grone1984positive}.  Here the
known object is the compression \(A^\top P A\), and the missing object is the
entire positive-definite fiber.  The selected completion is defined by AIRM
distance to a reference and is characterized globally as a horizontal
section of a metric submetry.

Best approximation under AIRM has been studied for geodesically convex,
unitarily invariant subsets through logarithmic and spectral reductions
\citep{bhatia2014approximation}, and for Riemannian geodesic submanifolds
\citep{lim2004best}.  Totally geodesic AIRM submanifolds and their
nearest-point projections have also been characterized directly inside the SPD
manifold \citep{tumpach2024totally}.  Our information sheet is a particular
congruence-transformed block submanifold of that kind.  The additional object
here is the transverse affine-compression fiber: we identify its unique
nearest point to every reference, prove that the compression is a metric
submetry, and derive universal radial-visible reduction.  Thus total geodesy
supports the construction but does not by itself supply the fiber-completion
or variational statements.

The distinction can be stated theorem by theorem.  Existing AIRM
approximation results supply general uniqueness or explicit projections for
specified convex or totally geodesic target sets.  In
\cref{thm:spd-compression}, the constrained set is instead the generally
noncompact fiber \(\{P:A^\top PA=R\}\).  Contraction of the compression gives
the lower bound, while a transverse block section attains it for every target.
The new package is the equality-attaining unique fiber lift, its coherent
horizontal foliation and submetry consequences, followed by the
moving-channel pullback metric of \cref{thm:stratified-line-element}.  We use
the classical approximation and totally geodesic theories as inputs to this
positioning, rather than claiming that AIRM nearest-point uniqueness by itself
is new.

Quotient geometries for fixed-rank positive-semidefinite matrices have also
been developed for matrix means and optimization
\citep{bonnabel2010fixedrank,vandereycken2013fixedrank,massart2020quotient}.
Those manifolds parameterize a rank-\(m\) PSD matrix.  Our stratum instead
fixes the rank of the displacement \(P-P_0\), which may be indefinite while
\(P\) remains positive definite, and its line element is the pullback of the
ambient AIRM at \(P_0+(P-P_0)\).  The resulting mismatch weight
\(R+R^{-1}-2I\) and its reference-valued rotation kernel are specific to this
completion geometry.
Viewed only as a set, each fixed-inertia component of this displacement
stratum is an open part of the standard fixed-rank symmetric matrix manifold.
The contribution is therefore not the existence of that manifold alone; it is
the global completion-descriptor diffeomorphism and the ambient-AIRM pullback,
whose mismatch weight and reference-valued rotation kernel do not arise from
the usual Euclidean fixed-rank metric.

\paragraph{Riemannian submersions and submetries.}
The differential geometry of Riemannian submersions starts from the
horizontal--vertical equations of O'Neill \citep{oneill1966submersion}.
Metric submetries provide the corresponding ball-surjectivity language;
Berestovskii and Guijarro \citep{berestovskii2000metric} characterize
Riemannian submersions through this metric property, and Guijarro and
Walschap \citep{guijarro2011submetries} delineate when submetries recover
smooth submersions.  Our abstract theorem uses this
classical metric structure and adds a variational characterization: exact
fiber distance is equivalent to lossless reduction for every monotone radial
visible objective.  The split-Hadamard result imposes global integrable
horizontal geometry to obtain coherent shortest sections, while the SPD
compression theorem verifies this structure and resolves the section in
closed form.

\paragraph{Information geometry and natural gradient.}
Natural gradient equips statistical model space with Fisher geometry
\citep{amari1998natural}.  Structured Fisher and Gauss--Newton approximations
motivate K-FAC, EK-FAC, and related methods
\citep{martens2015kfac,george2018ekfac,martens2020natural}; limitations of
identifying Fisher and Hessian geometry are discussed by Kunstner, Balles,
and Hennig \citep{kunstner2019limitations}.  Mirror descent also admits an
information-geometric interpretation \citep{raskutti2015information}.  Our
focus is the observation layer beneath a chosen training geometry: which
cotangent information is visible, which ambient completions are compatible,
and which completion is selected by a reference.

Sufficient statistics contract information-geometric tensors and preserve
them under stronger sufficiency conditions \citep{ay2015sufficient}.  That
work concerns statistical models and Fisher-type tensors.  Here the declared
object is a finite-dimensional positive cometric, the observation is the
linear compression \(A^\top PA\), and the main question is its
reference-dependent AIRM completion.  Statistical sufficiency motivates the
visibility language, while the completion and moving-channel line element
address a different geometric problem.
The AIRM covariance factor is also the Fisher geometry of the multivariate
normal family up to convention-dependent scaling
\citep{lovric2000multivariate}; this explains the Gaussian interpretation but
does not determine the reference-relative completion of an arbitrary declared
visible cometric.

\paragraph{Metric learning and relative-entropy projections.}
Information-theoretic metric learning uses LogDet or relative-entropy
projections under supervised constraints \citep{davis2007information}.  Our
Gaussian KL consequence is related, but the main object is the full fiber of
an information compression.  The metric-submetry theorem explains when every
radial visible decision, not only a particular divergence projection, reduces
without loss.  Classical information projections minimize divergence under
linear or convex constraints
\citep{csiszar1975idivergence,csiszar2003projections}.  Our Gaussian KL
corollary is a finite-dimensional instance with a closed-form fiber
projection; the AIRM shortest-completion and stratified line-element claims
are separate from the divergence-projection theory.

\paragraph{Adaptive and structured preconditioning.}
AdaGrad, Adam, Shampoo, and related methods impose diagonal, block, matrix, or
tensor structure on positive preconditioners
\citep{duchi2011adagrad,kingma2014adam,gupta2018shampoo,morwani2024shampoo}.
Optimal diagonal preconditioning and matrix-scaling viewpoints study related
coordinate restrictions \citep{qu2025optimaldiagonal}.  We formulate a
different certificate: whether a structured positive family can realize a
declared visible target.  The answer is a relative-interior condition in a
linear conic image, which distinguishes strict representability, boundary
failure, and closed-cone infeasibility.

\paragraph{Convex certificates.}
The structured certificates use finite-dimensional convex separation,
relative interiors, and semidefinite images
\citep{rockafellar1970convex,vandenberghe1996semidefinite,boyd2004convex}.
These tools enter only after the full-SPD information map has been resolved;
they quantify what is lost when the completion is restricted to an
implementable optimizer family.

\paragraph{Subspace perturbation and matrix concentration.}
When the visible subspace is estimated from a covariance or Gram matrix, its
recovery requires an eigengap and a principal-angle perturbation bound
\citep{yu2015davis}.  Covariance concentration then controls the ambient
operator error \citep{tropp2015matrix}.  We combine these ingredients with
the exact stratified line element to obtain a completed-geometry error bound
that separates visible SPD error from subspace rotation.
